\section{Bounded K-Level Synthesis Support Summary}

\label{sec:oc133-bounded-klevel-synthesis-support-summary}

This appendix replaces raw generated K-level status dossiers with a reader-facing synthesis summary. The machine-readable evidence package remains the place for exhaustive rows, hashes, and verdict fields. The public monograph keeps the scientific question visible: what structural role does each K-level play, what evidence class supports it, and what would reopen the claim?


\subsection{K0: proto-continuum and zero-continuumness boundary}

K0 is included as a declared witness-taxonomy level, not as an unbounded domain-completion claim. Its current release role is proto-continuum and zero-continuumness boundary. The supporting evidence class is K0 proof route and finite semantic witness. A reviewer should read the level as a scoped part of the model grammar and should not infer that every possible phenomenon at this level has already been explained or numerically predicted.


The K0 reopening condition is local. If the claimed level distinction can be projected away without changing any declared witness, proof route, replay row, comparator boundary, or phenomenon model card, then the level-specific public wording is too strong and must be repaired or demoted.


\subsection{K1: low-level structure, separation, and early organization}

K1 is included as a declared witness-taxonomy level, not as an unbounded domain-completion claim. Its current release role is low-level structure, separation, and early organization. The supporting evidence class is typed model definitions, boundary predicates, and representative finite checks. A reviewer should read the level as a scoped part of the model grammar and should not infer that every possible phenomenon at this level has already been explained or numerically predicted.


The K1 reopening condition is local. If the claimed level distinction can be projected away without changing any declared witness, proof route, replay row, comparator boundary, or phenomenon model card, then the level-specific public wording is too strong and must be repaired or demoted.


\subsection{K2: low-level structure, separation, and early organization}

K2 is included as a declared witness-taxonomy level, not as an unbounded domain-completion claim. Its current release role is low-level structure, separation, and early organization. The supporting evidence class is typed model definitions, boundary predicates, and representative finite checks. A reviewer should read the level as a scoped part of the model grammar and should not infer that every possible phenomenon at this level has already been explained or numerically predicted.


The K2 reopening condition is local. If the claimed level distinction can be projected away without changing any declared witness, proof route, replay row, comparator boundary, or phenomenon model card, then the level-specific public wording is too strong and must be repaired or demoted.


\subsection{K3: organized physical, chemical, and excitable structures}

K3 is included as a declared witness-taxonomy level, not as an unbounded domain-completion claim. Its current release role is organized physical, chemical, and excitable structures. The supporting evidence class is bounded reconstruction rows and formal boundary conditions. A reviewer should read the level as a scoped part of the model grammar and should not infer that every possible phenomenon at this level has already been explained or numerically predicted.


The K3 reopening condition is local. If the claimed level distinction can be projected away without changing any declared witness, proof route, replay row, comparator boundary, or phenomenon model card, then the level-specific public wording is too strong and must be repaired or demoted.


\subsection{K4: organized physical, chemical, and excitable structures}

K4 is included as a declared witness-taxonomy level, not as an unbounded domain-completion claim. Its current release role is organized physical, chemical, and excitable structures. The supporting evidence class is bounded reconstruction rows and formal boundary conditions. A reviewer should read the level as a scoped part of the model grammar and should not infer that every possible phenomenon at this level has already been explained or numerically predicted.


The K4 reopening condition is local. If the claimed level distinction can be projected away without changing any declared witness, proof route, replay row, comparator boundary, or phenomenon model card, then the level-specific public wording is too strong and must be repaired or demoted.


\subsection{K5: organized physical, chemical, and excitable structures}

K5 is included as a declared witness-taxonomy level, not as an unbounded domain-completion claim. Its current release role is organized physical, chemical, and excitable structures. The supporting evidence class is bounded reconstruction rows and formal boundary conditions. A reviewer should read the level as a scoped part of the model grammar and should not infer that every possible phenomenon at this level has already been explained or numerically predicted.


The K5 reopening condition is local. If the claimed level distinction can be projected away without changing any declared witness, proof route, replay row, comparator boundary, or phenomenon model card, then the level-specific public wording is too strong and must be repaired or demoted.


\subsection{K6: cognitive, social, institutional, and systems-scale organization}

K6 is included as a declared witness-taxonomy level, not as an unbounded domain-completion claim. Its current release role is cognitive, social, institutional, and systems-scale organization. The supporting evidence class is model-card routes, replay QA rows, and comparator boundaries. A reviewer should read the level as a scoped part of the model grammar and should not infer that every possible phenomenon at this level has already been explained or numerically predicted.


The K6 reopening condition is local. If the claimed level distinction can be projected away without changing any declared witness, proof route, replay row, comparator boundary, or phenomenon model card, then the level-specific public wording is too strong and must be repaired or demoted.


\subsection{K7: cognitive, social, institutional, and systems-scale organization}

K7 is included as a declared witness-taxonomy level, not as an unbounded domain-completion claim. Its current release role is cognitive, social, institutional, and systems-scale organization. The supporting evidence class is model-card routes, replay QA rows, and comparator boundaries. A reviewer should read the level as a scoped part of the model grammar and should not infer that every possible phenomenon at this level has already been explained or numerically predicted.


The K7 reopening condition is local. If the claimed level distinction can be projected away without changing any declared witness, proof route, replay row, comparator boundary, or phenomenon model card, then the level-specific public wording is too strong and must be repaired or demoted.


\subsection{K8: cognitive, social, institutional, and systems-scale organization}

K8 is included as a declared witness-taxonomy level, not as an unbounded domain-completion claim. Its current release role is cognitive, social, institutional, and systems-scale organization. The supporting evidence class is model-card routes, replay QA rows, and comparator boundaries. A reviewer should read the level as a scoped part of the model grammar and should not infer that every possible phenomenon at this level has already been explained or numerically predicted.


The K8 reopening condition is local. If the claimed level distinction can be projected away without changing any declared witness, proof route, replay row, comparator boundary, or phenomenon model card, then the level-specific public wording is too strong and must be repaired or demoted.


\subsection{K9: theory, meta-theory, and cross-model coherence under declared assumptions}

K9 is included as a declared witness-taxonomy level, not as an unbounded domain-completion claim. Its current release role is theory, meta-theory, and cross-model coherence under declared assumptions. The supporting evidence class is proof/evidence governance, comparator positioning, and adversarial reopening rules. A reviewer should read the level as a scoped part of the model grammar and should not infer that every possible phenomenon at this level has already been explained or numerically predicted.


The K9 reopening condition is local. If the claimed level distinction can be projected away without changing any declared witness, proof route, replay row, comparator boundary, or phenomenon model card, then the level-specific public wording is too strong and must be repaired or demoted.


\subsection{K10: theory, meta-theory, and cross-model coherence under declared assumptions}

K10 is included as a declared witness-taxonomy level, not as an unbounded domain-completion claim. Its current release role is theory, meta-theory, and cross-model coherence under declared assumptions. The supporting evidence class is proof/evidence governance, comparator positioning, and adversarial reopening rules. A reviewer should read the level as a scoped part of the model grammar and should not infer that every possible phenomenon at this level has already been explained or numerically predicted.


The K10 reopening condition is local. If the claimed level distinction can be projected away without changing any declared witness, proof route, replay row, comparator boundary, or phenomenon model card, then the level-specific public wording is too strong and must be repaired or demoted.


\subsection{K11: theory, meta-theory, and cross-model coherence under declared assumptions}

K11 is included as a declared witness-taxonomy level, not as an unbounded domain-completion claim. Its current release role is theory, meta-theory, and cross-model coherence under declared assumptions. The supporting evidence class is proof/evidence governance, comparator positioning, and adversarial reopening rules. A reviewer should read the level as a scoped part of the model grammar and should not infer that every possible phenomenon at this level has already been explained or numerically predicted.


The K11 reopening condition is local. If the claimed level distinction can be projected away without changing any declared witness, proof route, replay row, comparator boundary, or phenomenon model card, then the level-specific public wording is too strong and must be repaired or demoted.


\subsection{K12: theory, meta-theory, and cross-model coherence under declared assumptions}

K12 is included as a declared witness-taxonomy level, not as an unbounded domain-completion claim. Its current release role is theory, meta-theory, and cross-model coherence under declared assumptions. The supporting evidence class is proof/evidence governance, comparator positioning, and adversarial reopening rules. A reviewer should read the level as a scoped part of the model grammar and should not infer that every possible phenomenon at this level has already been explained or numerically predicted.


The K12 reopening condition is local. If the claimed level distinction can be projected away without changing any declared witness, proof route, replay row, comparator boundary, or phenomenon model card, then the level-specific public wording is too strong and must be repaired or demoted.


This summary preserves the scientific content needed for the monograph while preventing generated status fields from becoming public prose. Exhaustive K-level rows remain available in the public evidence package for reviewers who need to inspect them mechanically.
